\documentclass[11pt,a4paper]{article}
\usepackage[utf8]{inputenc}
\usepackage[vietnamese,english]{babel}
\usepackage{amsmath,amsfonts,amssymb,amsthm}
\usepackage{geometry}
\usepackage{xcolor}
\usepackage{titlesec}
\usepackage{hyperref}

% Cấu hình lề trang tiêu chuẩn cho báo cáo học thuật
\geometry{top=2cm,bottom=2cm,left=2cm,right=2cm}

% Tối ưu hiển thị cấu trúc tiêu đề các phần
\titleformat{\section}{\color{blue!70!black}\bfseries\large}{\thesection}{1em}{}[\titlerule]
\titleformat{\subsection}{\color{blue!50!black}\bfseries\mdseries}{\thesubsection}{1em}{}

\begin{document}

% =========================================================================
% PHẦN 1: TOÀN BỘ NỘI DUNG TIẾNG VIỆT
% =========================================================================
\selectlanguage{vietnamese}

\begin{center}
    \textbf{\LARGE BÀI LÀM TỰ LUẬN CHUYÊN ĐỀ ỨNG DỤNG AI} \\
    \vspace{0.2cm}
    \textbf{\large Chủ đề: Multiple Drone Attack in Urban City} \\
    \textit{(Tối ưu hóa chiến lược phòng thủ chống tấn công đa mục tiêu bằng thiết bị bay không người lái trong môi trường đô thị)}
\end{center}

\vspace{0.5cm}

\section{PHẦN 1: ĐẶT VẤN ĐỀ \& MỤC TIÊU (Chiếm 15\% dung lượng)}

\begin{itemize}
    \item \textbf{Bối cảnh thực tế:} Trong bối cảnh an ninh hiện đại, sự gia tăng của các thiết bị bay không người lái (drone) giá rẻ tạo ra mối đe dọa cực lớn đối với an ninh đô thị \cite{mutzari2025defending}. Các tác nhân độc hại có thể phối hợp điều khiển nhiều drone tấn công đồng thời (\textit{Multi-drone attacks}) nhắm vào các hạ tầng trọng yếu như cơ quan chính phủ, nhà máy điện, hệ thống nước và bệnh viện \cite{mutzari2025defending}. Môi trường đô thị với đặc thù nhiều tòa nhà cao tầng, vật cản và nguy cơ nhiễu tín hiệu là một thách thức không thể vượt qua đối với các hệ thống phòng thủ mặt đất truyền thống \cite{mutzari2025defending}. Do drone có lợi thế giá rẻ, cơ động cao, thành phố cần được bảo vệ bằng một phi đội drone đánh chặn (\textit{Defense drones}) hoạt động trên không \cite{mutzari2025defending}.
    \item \textbf{Mục tiêu cốt lõi:} Đồ án này không tiếp cận theo tư duy thụ động, bám đuổi vu vơ ``thấy đâu bắt đó'' của các hệ thống cũ \cite{mutzari2025defending}. Thay vào đó, mục tiêu là ứng dụng AI chiến lược nhằm \underline{tối thiểu hóa thiệt hại về người và tài sản} (\textit{Minimize Damage}) cho toàn thành phố thông qua việc tối ưu hóa phân bổ nguồn lực \cite{mutzari2025defending}. Đồ án chấp nhận một thực tế khách quan rằng tài nguyên phòng thủ của ta luôn bị giới hạn so với số lượng drone của kẻ địch \cite{mutzari2025defending}.
\end{itemize}

\section{PHẦN 2: MÔ HÌNH HÓA TOÁN HỌC BÀI TOÁN (Chiếm 25\% dung lượng)}

Để giải quyết bài toán điều phối chiến lược, hệ thống áp dụng Lý thuyết trò chơi an ninh (\textit{Security Games}) với các định nghĩa toán học chặt chẽ \cite{mutzari2025defending}:
\begin{itemize}
    \item \textbf{Mô hình đồ thị:} Không gian đô thị được mô hình hóa dưới dạng một đồ thị tổng quát $G=(V,E)$ \cite{mutzari2025defending}. Trong đó, tập hợp các nút $V$ đại diện cho các khu vực hoặc mục tiêu cần bảo vệ trên mặt đất, và tập hợp các cạnh $E$ đại diện cho các tuyến đường bay kết nối giữa các khu vực \cite{mutzari2025defending}.
    \item \textbf{Mô hình Trò chơi SSG tuần tự:} Trận chiến được thiết lập dưới dạng một Trò chơi an ninh Stackelberg tuần tự (\textit{Sequential Stackelberg Security Game - SSSG}) \cite{mutzari2025defending}. Bên phòng thủ (\textit{Defender} - hệ thống AI đô thị) đóng vai trò người đi trước, cam kết một chiến lược phối hợp ngẫu nhiên (\textit{Mixed strategy}) để phân bổ các drone đánh chặn \cite{mutzari2025defending}. Kẻ tấn công (\textit{Attacker}) đóng vai trò người đi sau, giám sát chiến lược phân bổ này và đưa ra phản ứng tối ưu nhất (\textit{Best response}) nhằm tối đa hóa thiệt hại cho ta \cite{mutzari2025defending}.
    \item \textbf{Các ràng buộc thực tế:} Hệ thống đưa vào các tham số vật lý bao gồm: Dung lượng pin (Battery $B$) quyết định số bước di chuyển tối đa của drone (mỗi cạnh di chuyển hoặc loitering tiêu tốn đúng 1 đơn vị pin) và Tải trọng (Payload $P$) quy định số lượng mục tiêu tối đa mà một drone tấn công có thể phá hủy \cite{mutzari2025defending}.
    \item \textbf{Hàm mục tiêu (\textit{General-sum Utility}):} Trò chơi được thiết kế theo dạng tổng quát (\textit{General-sum}), nghĩa là lợi ích hai bên không đối nghịch hoàn toàn \cite{mutzari2025defending}. Trọng số phần thưởng của kẻ tấn công ($R^a$) và hình phạt của bên phòng thủ ($P^d$) được thiết lập dựa trên độ quan trọng của hạ tầng đô thị được khảo sát bởi chuyên gia bảo mật \cite{mutzari2025defending}.
\end{itemize}

\section{PHẦN 3: PHÂN TÍCH GIẢI THUẬT CỐT LÕI S2D2 (Chiếm 40\% dung lượng)}

Do không gian trạng thái của bài toán tuần tự tỷ lệ thuận theo hàm mũ với quy mô đô thị, việc tìm trạng thái cân bằng Strong Stackelberg Equilibrium (SSE) chính xác là một bài toán NP-hard \cite{korzhyk2011security}. Đồ án đề xuất giải thuật \textbf{S2D2 (Sequential Stackelberg Drone Defense)} nhằm tìm ra nghiệm xấp xỉ tối ưu thông qua 3 bước xử lý \cite{mutzari2025defending}:

\subsection*{Sơ đồ luồng xử lý hệ thống (Pipeline Architecture)}
\begin{center}
    \small $\text{[Bản đồ Đô thị } G=(V,E)\text{]} \xrightarrow{\text{Bước 1: Coarsening}} \text{[Cụm Khu lân cận]} \xrightarrow{\text{Bước 2: Giải Subgame DFS}} \text{[Tối ưu đơn drone]} \xrightarrow{\text{Bước 3: Giải Meta-Game}} \text{[Ma trận phân bổ MILP]}$
\end{center}

\subsection*{Chi tiết 3 bước xử lý của thuật toán:}
\begin{enumerate}
    \item \textbf{Bước 1: Thu nhỏ đồ thị (Graph Coarsening - $\delta$-coarsening):} Bản đồ thành phố thực tế vô cùng khổng lồ với hàng trăm nghìn nút, khiến AI không thể tính toán trực tiếp \cite{mutzari2025defending}. S2D2 áp dụng phân cụm đồ thị thành các cụm nhân tạo gọi là ``Khu lân cận'' (\textit{Neighborhoods}) dựa trên tham số granularity $\delta$ \cite{mutzari2025defending}. Thuật toán bỏ qua các nút có phần thưởng nhỏ hơn $\delta$ để tập trung cô lập và bảo vệ các cụm mục tiêu có giá trị cao \cite{mutzari2025defending}.
    \item \textbf{Bước 2: Giải bài toán đơn mục tiêu (Single-Attacker Single-Defender Subgame):} Trong phạm vi từng khu lân cận, hệ thống tính toán chiến lược đối phó tối ưu giữa 1 drone phòng thủ và 1 drone tấn công dưới một xác suất hiện diện $\lambda$ giả định của defender \cite{korzhyk2010complexity}. Hệ thống áp dụng kỹ thuật duyệt cây DFS sâu kết hợp loại bỏ các chiến lược bị thống trị (\textit{Dominated strategies}) để thu hẹp không gian tìm kiếm \cite{mutzari2025defending}.
    \item \textbf{Bước 3: Giải trò chơi tổng thể (Solving the Meta-Game):} Sau khi có hàm lợi ích từ các khu lân cận, S2D2 quy đổi toàn bộ bài toán điều phối phi đội drone đa mục tiêu thành một bài toán Quy hoạch tuyến tính nguyên hỗn hợp (\textbf{MILP} - \textit{Mixed Integer Linear Program}) \cite{mutzari2025defending}. Bằng cách áp dụng xấp xỉ tuyến tính từng đoạn theo xác suất $\lambda$, mô hình MILP sẽ được giải để đưa ra ma trận phân bổ hỗn hợp ngẫu nhiên, xác định vị trí xuất phát tối ưu nhất cho từng drone đánh chặn \cite{mutzari2025defending}.
\end{enumerate}

\section{PHẦN 4: KẾT QUẢ THỰC NGHIỆM \& TƯ DUY PHẢN BIỆN (Chiếm 20\% dung lượng)}

\begin{itemize}
    \item \textbf{Chứng minh thực nghiệm:} Hệ thống đã được kiểm thử thực tế trên tập dữ liệu nâng cao trích xuất thông qua thư viện \texttt{OSMnx} \cite{boeing2017osmnx} của 80 thành phố lớn trên thế giới (quy mô lên tới 277,000 nút) \cite{mutzari2025defending}. Kết quả định lượng cho thấy thuật toán S2D2 giúp giảm thiểu tối đa thiệt hại hạ tầng, \textbf{tăng hiệu quả bảo vệ lên trung bình 2.84 lần đối với dữ liệu mô phỏng} và \textbf{từ 34\% đến 41\% đối với dữ liệu cấu hình thực tế từ chuyên gia} khi so sánh với các giải thuật bầy đàn truyền thống \cite{mutzari2025defending}. Về mặt thời gian xử lý, hệ thống chỉ mất vài phút để giải toán, hoàn toàn khả thi để triển khai thực tế \cite{mutzari2025defending}.
    \item \textbf{Tư duy phản biện (Hạn chế của hệ thống):} Mô hình hiện tại giả định kẻ tấn công di chuyển theo chiến lược tĩnh (\textit{Open-loop}) do không biết vị trí drone phòng thủ trước khi khai hỏa \cite{mutzari2025defending}. Tuy nhiên, các thử nghiệm làm nhiễu thông tin phần thưởng từ 0\% đến 100\% chứng minh chiến lược S2D2 vẫn giữ được xu hướng ổn định và có tính bền vững (\textit{Robustness}) cao \cite{mutzari2025defending}.
    \item \textbf{Hướng phát triển tương lai:} Đồ án định hướng mở rộng mô hình trò chơi với các drone không đồng nhất (khác biệt vận tốc, có camera nhiệt quét đêm) \cite{mutzari2025defending}. Đồng thời, tích hợp hệ thống toán chiến lược này với các module thị giác máy tính nhận diện mục tiêu thời gian thực như YOLO và DeepSORT để tạo nên một lá chắn đô thị khép kín \cite{mutzari2025defending}.
\end{itemize}


% =========================================================================
% PHẦN 2: TOÀN BỘ NỘI DUNG TIẾNG ANH (CHUYỂN SANG TRANG MỚI)
% =========================================================================
\newpage
\selectlanguage{english}

\begin{center}
    \textbf{\LARGE EXECUTIVE ARTIFICIAL INTELLIGENCE ESSAY} \\
    \vspace{0.2cm}
    \textbf{\large Topic: Multiple Drone Attack in Urban City} \\
    \textit{(Strategic Defense Optimization Against Multi-Target UAV Attacks in Urban Environments)}
\end{center}

\vspace{0.5cm}

\section{PART 1: PROBLEM STATEMENT \& OBJECTIVE (15\% Weight)}

\begin{itemize}
    \item \textbf{Real-World Context:} In modern security paradigms, the proliferation of low-cost Unmanned Aerial Vehicles (drones) presents an asymmetric threat to urban security \cite{mutzari2025defending}. Malicious actors can deploy coordinated multi-drone attacks targeting critical infrastructure such as government buildings, power grids, water supply networks, and hospitals \cite{mutzari2025defending}. Urban topographies, characterized by high-rise buildings, multi-layered obstacles, and signal interference, completely neutralize traditional ground-based radar and defensive hardware \cite{mutzari2025defending}. Consequently, deploying a highly mobile, aerial interceptor squadron of defense drones is imperative for urban protection \cite{mutzari2025defending}.
    \item \textbf{Core Objective:} This project rejects the passive, reactive ``catch-as-catch-can'' approach of legacy tracking systems \cite{mutzari2025defending}. Instead, the objective is to implement strategic AI frameworks designed to \underline{minimize total systemic damage} to human lives and physical assets by optimizing defense resource allocation \cite{mutzari2025defending}. The framework models the problem under the operational constraint that defensive assets are inherently limited compared to the adversary's offensive capabilities \cite{mutzari2025defending}.
\end{itemize}

\section{PART 2: MATHEMATICAL MODELING OF THE GAME (25\% Weight)}

To resolve the tactical deployment problem, the system models the engagement using game-theoretic principles through Security Games \cite{mutzari2025defending}:
\begin{itemize}
    \item \textbf{Graph Modeling:} The urban airspace is abstracted as a general undirected graph $G=(V,E)$ \cite{mutzari2025defending}. The vertex set $V$ defines targeted ground zones or critical infrastructure sectors requiring physical protection, while the edge set $E$ maps the interconnected aerial flight corridors connecting these sectors \cite{mutzari2025defending}.
    \item \textbf{Sequential SSSG Model:} The tactical encounter is formulated as a \textbf{Sequential Stackelberg Security Game (SSSG)} \cite{mutzari2025defending}. The Defender (the automated urban AI backend) acts as the leader, committing to a randomized \textbf{Mixed Strategy} profile to allocate aerial interceptors across sectors \cite{mutzari2025defending}. The Attacker (the adversary) acts as the follower, monitoring the defender's commitment map to deploy a mathematically calculated \textbf{Best Response} flight path that maximizes systemic damage \cite{mutzari2025defending}.
    \item \textbf{Hardware Constraints:} The environment enforces strict physical parameters, including: Battery capacity ($B$) which restricts the maximal flight path length of the drones (where each edge traversal or hovering action consumes exactly 1 battery unit) and Payload capacity ($P$) which dictates the maximum number of targets an attacker drone can neutralize \cite{mutzari2025defending}.
    \item \textbf{Utility Formulation (General-Sum Game):} The environment is configured as a \textbf{General-Sum game}, acknowledging that the players' objectives are not diametrically opposed \cite{mutzari2025defending}. Attacker reward metrics ($R^a$) and defender penalty weights ($P^d$) are calibrated independently based on priority evaluations by security experts \cite{mutzari2025defending}.
\end{itemize}

\section{PART 3: S2D2 CORE ALGORITHMIC ANALYSIS (40\% Weight)}

Because the extensive-form state space of a sequential graph game scales exponentially with urban dimensions, computing an exact Strong Stackelberg Equilibrium (SSE) is an NP-hard problem \cite{korzhyk2011security}. This project proposes the \textbf{S2D2 (Sequential Stackelberg Drone Defense)} paradigm to isolate approximate equilibrium strategies through a three-step hierarchical framework \cite{mutzari2025defending}:

\subsection*{Algorithmic Architecture Pipeline (Pipeline Architecture)}
\begin{center}
    \small $\text{[Urban Graph } G=(V,E)\text{]} \xrightarrow{\text{Step 1: Graph Coarsening}} \text{[Neighborhood Clusters]} \xrightarrow{\text{Step 2: Subgame Engine (DFS)}} \text{[Single-Drone Subgame]} \xrightarrow{\text{Step 3: Meta-Game Solver}} \text{[MILP Matrix]}$
\end{center}

\subsection*{Detailed Algorithmic Execution Steps:}
\begin{enumerate}
    \item \textbf{Step 1: Graph Coarsening ($\delta$-coarsening):} Real-world urban maps scale to hundreds of thousands of nodes, rendering direct global matrix computations impossible \cite{mutzari2025defending}. S2D2 applies an adaptive clustering routine to partition the topology into isolated artificial ``Neighborhoods'' based on a granularity parameter $\delta$ \cite{mutzari2025defending}. The routine drops low-value sectors where rewards fall below $\delta$, effectively isolating high-value targets into tightly bound sub-graphs \cite{mutzari2025defending}.
    \item \textbf{Step 2: Single-Attacker Single-Defender Subgame Engine:} Within each isolated neighborhood cluster, the framework evaluates a localized execution tree between 1 defender and 1 attacker drone under a probabilistic defender presence coefficient $\lambda$ \cite{korzhyk2010complexity}. The engine applies deep Depth-First Search (DFS) routines coupled with active pruning of mathematically \textbf{Dominated Strategies} to contract the search space and accelerate convergence \cite{mutzari2025defending}.
    \item \textbf{Step 3: Global Meta-Game Solver:} After extracting parameterized utility functions across neighborhoods, S2D2 reduces the decentralized multi-drone allocation problem back into a static, multi-resource assignment model resolved via a \textbf{Mixed Integer Linear Program (MILP)} \cite{mutzari2025defending}. By conducting a piece-wise linear approximation over the variable probability vector $\lambda$, the global MILP is compiled and resolved to output the definitive mixed strategy deployment matrix for the interceptor fleet \cite{mutzari2025defending}.
\end{enumerate}

\section{PART 4: EMPIRICAL VALIDATION \& CRITICAL DISCUSSION (20\% Weight)}

\begin{itemize}
    \item \textbf{Quantitative Empirical Validation:} The defensive architecture was validated using real-world open street networks extracted via the \texttt{OSMnx} library \cite{boeing2017osmnx} across 80 international metropolises, scaling up to 277,000 nodes \cite{mutzari2025defending}. The quantitative metrics reveal that the S2D2 engine dramatically limits physical destruction, achieving an average \textbf{2.84x defender utility improvement in simulated configurations} and a \textbf{34\% to 41\% utility preservation gain on maps annotated by security experts} when compared to baseline greedy drone-swarm architectures \cite{mutzari2025defending}. In terms of operational execution speeds (\textit{Runtime Scalability}), the solver successfully converges within a few minutes for standard city sizes, proving its viability for real-world strategic deployment \cite{mutzari2025defending}.
    \item \textbf{Systemic Critique (Limitations):} The current tactical model assumes an Open-loop adversary (modeling an attacker blind to active interceptor vectors during flight) \cite{mutzari2025defending}. However, rigorous reward perturbation assessments scaling from 0\% up to 100\% demonstrate that the deployment matrices generated by S2D2 exhibit extraordinary structural robustness against adversarial behavioral adaptations \cite{mutzari2025defending}.
    \item \textbf{Future Research Horizons:} Future milestones aim to expand the mathematical framework to support heterogeneous drone parameters (variable velocities and thermal scanning hardware) \cite{mutzari2025defending}. Furthermore, we plan to fuse this game-theoretic strategic layer directly with live real-time computer vision networks (YOLO object detection + DeepSORT multi-object tracking) to establish an impenetrable, fully autonomous urban iron-dome \cite{mutzari2025defending}.
\end{itemize}

\newpage
% Tự động nạp dữ liệu từ tệp references.bib để in danh mục tài liệu tham khảo chung ở cuối bài
\bibliographystyle{plain}
\bibliography{references}

\end{document}